% ============ 第3章 参考文献(原书 pp.305-309) ============
\section*{参考文献}\addcontentsline{toc}{section}{参考文献}
\begin{thebibliography}{99}
% 条目照录原文(书名/作者不翻译).页码范围用 en dash.
\bibitem[1]{ch3:1} ブルバキ(倉田令二朗訳),数学原論 スペクトル論 I, 東京図書, 東京, 1971.
\bibitem[2]{ch3:2} 日本数学会 編,関数環とその関連分野(特集),数学 28(1976), no.1.
\bibitem[3]{ch3:3} 荷見守助,調和解析学,槙書店,東京, 1978.
\begin{sloppypar}
\bibitem[4]{ch3:4} J. Bergh and J. Löfström, Interpolation Spaces: An introduction, Springer-Verlag, Berlin-Heidelberg-New York, 1976.
\end{sloppypar}
\bibitem[5]{ch3:5} J. Duoandikoetxea, Fourier analysis, Graduate Studies in Mathematics, A.M.S, 2001.
\bibitem[6]{ch3:6} C. C. Graham, O. C. McGehee, Essays in commutative harmonic analysis, Springer-Verlag, New York-Heidelberg-Berlin, 1979.
\bibitem[7]{ch3:7} J. P. Kahane, Séries de Fourier absolument convergentes, Springer-Verlag, Berlin, 1970.
\bibitem[8]{ch3:8} Y. Katznelson, An introduction to harmonic analysis, Dover Publications, New York, 1976.
\bibitem[9]{ch3:9} R. Larsen, An Introduction to the Theory of Multipliers, Springer-Verlag, Berlin-Heidelberg-New York, 1971.
\bibitem[10]{ch3:10} L. H. Loomis, An introduction to abstract harmonic analysis, Princeton, Van Nostrand, 1953.
\bibitem[11]{ch3:11} W. Rudin, Fourier analysis on groups, Interscience, New York-London, 1962.
\bibitem[12]{ch3:12} J. L. Taylor, Measure algebras, CBMS Conference Report, No. 16, A.M.S, 1973.
\bibitem[13]{ch3:13} G. Brown, Construction of Fourier multipliers, Bull. Austral. Math. Soc. 16(1977), 463--472.
\bibitem[14]{ch3:14} A. P. Calderón and A. Zygmund, Algebras of certain singular operators, Amer. J. Math. 78(1956), 310--320.
\bibitem[15]{ch3:15} Y. Domar, On the spectral synthesis problem for $(n-1)$-dimensional subsets of $R^{n}$, $n \geq 2$, Ark. Mat. 9(1971), 23--37.
\bibitem[16]{ch3:16} R. Doss, Some inclusions in multipliers, Pacific J. Math. 32(1970), 643--646.
\bibitem[17]{ch3:17} K. de Leeuw, On $L_{p}$ multipliers, Ann. of Math. (2) 81(1965), 364--379.
\bibitem[18]{ch3:18} A. Figà-Talamanca and G. I. Gaudry, Multipliers of $L^{p}$ which vanish at infinity, J. Functional Analysis 7(1971), 475--486.
\bibitem[19]{ch3:19} O. Hatori, Functions which operate on algebras of Fourier multipliers, Tohoku Math. J. (2) 47(1995), no. 1, 105--115.
\bibitem[20]{ch3:20} O. Hatori, A characterization of lacunary set and spectral properties of Fourier multipliers, Function spaces (Edwardsville, IL, 1994), 183--203, Lecture Notes in Pure and Appl. Math., 172, Dekker, New York, 1995.
\bibitem[21]{ch3:21} C. S. Herz, Spectral synthesis for the circle, Ann. of Math. 68(1958), 709--712.
\bibitem[22]{ch3:22} L. Hörmander, Estimates for translation invariant operators in $L^{p}$ spaces, Acta Math. 104(1960), 93--140.
\bibitem[23]{ch3:23} S. Igari, スペクトル合成について, 数学 18(1966/1967), 215--222.
\bibitem[24]{ch3:24} S. Igari, Functions of $L^{p}$ multipliers, Tôhoku Math. J. 21(1969), 304--320.
\bibitem[25]{ch3:25} S. Igari, Functions of $L^{p}$-multipliers II, Tôhoku Math. J. (2) 26(1974), 555--561.
\bibitem[26]{ch3:26} S. Igari, E. Sato, Operating functions on Fourier multipliers, Tohoku Math. J. (2) 46(1994), no. 3, 357--366.
\bibitem[27]{ch3:27} M. Jr. Jodeit, Restrictions and extensions of Fourier multipliers, Studia Math. 34(1970), 215--226.
\bibitem[28]{ch3:28} M. Kaneko and E. Sato, Notes on transference of continuity from maximal Fourier multiplier operators on $R^{n}$ to those on $T^{n}$, Interdiscip. Inform. Sci. 4(1998), no. 1, 97--107.
\bibitem[29]{ch3:29} Y. Kanjin, A convolution measure algebra on the unit disc, Tôhoku Math. J. 28(1976), 105--115.
\bibitem[30]{ch3:30} Y. Kanjin, On algebras with convolution structures for Laguerre polynomials, Trans. Amer. Math. Soc. 295(1986), no. 2, 783--794.
\bibitem[31]{ch3:31} Y. Katznelson, Sur les fonctions opérant sur l'algèbre des séries de Fourier absolument convergentes, C. R. Acad. Sci. Paris 247(1958), 404--406.
\bibitem[32]{ch3:32} R. Kaufman, M-sets and distributions, Astérisque 5(1973), 225--230.
\bibitem[33]{ch3:33} T. Körner, A pseudofunction on a Helson set. I, Astérisque 5(1973), 3--224.
\bibitem[34]{ch3:34} P. Malliavin, Impossibilité de la synthèse spectrale sur les groupes abéliens non compacts, Publ. Math. Inst. Hautes Études Sci. 1959, 85--92.
\bibitem[35]{ch3:35} D. Müller, On the spectral synthesis problem for hypersurfaces of $\mathbf{R}^N$, J. Funct. Anal. 47(1982), no. 2, 247--280.
\bibitem[36]{ch3:36} F. Parreau, Measures with real spectra, Invent. Math. 98(1989), no. 2, 311--330.
\bibitem[37]{ch3:37} T. V. Pedersen, Some properties of the Pisier algebra, Math. Proc. Cambridge Philos. Soc. 128(2000), 343--354.
\bibitem[38]{ch3:38} G. Pisier, A remarkable homogeneous Banach algebra, Israel J. Math. 34(1979), no. 1-2, 38--44.
\bibitem[39]{ch3:39} I. Richards, On Malliavin's counterexample to spectral synthesis, Bull. Amer. Math. Soc. 72(1966), 698--700.
\bibitem[40]{ch3:40} D. Rider, Transformations of Fourier coefficients, Pacific J. Math. 19(1966), 347--355.
\bibitem[41]{ch3:41} W. Rudin, Some theorems on Fourier coefficients, Proc. Amer. Math. Soc. 10(1959), 855--859.
\bibitem[42]{ch3:42} W. Rudin, Fourier-Stieltjes transforms of measures on independent sets, Bull. Amer. Math. Soc. 66(1960), 199--202.
\bibitem[43]{ch3:43} W. Rudin, A strong converse of the Wiener--Lévy theorem, Canad. J. Math. 14(1962) 694--701.
\bibitem[44]{ch3:44} S. Saeki, Spectral synthesis for the Kronecker sets, J. Math. Soc. Japan 21(1969), 549--563.
\bibitem[45]{ch3:45} S. Saeki, Helson sets which disobey spectral synthesis, Proc. Amer. Math. Soc. 47(1975), 371--377.
\bibitem[46]{ch3:46} S. Saeki, Tensor products in Harmonic Analysis, KSU Lecture Note.
\bibitem[47]{ch3:47} P. Sarnak, Spectra of singular measures as multipliers on $L_{p}$, J. Functional Analysis 37(1980), 302--317.
\bibitem[48]{ch3:48} E. Sato, Spectra of measures as $L_{p}$ multipliers, Tohoku Math. J. (2) 37(1985), no. 3, 333--342.
\bibitem[49]{ch3:49} L. Schwartz, Sur une propriété de synthèse spectrale dans les groupes non compacts, C. R. Acad. Sci. Paris 227(1948), 424--426.
\bibitem[50]{ch3:50} N. Th. Varopoulos, Sur les ensembles parfaits et les séries trigonométriques, C. R. Acad. Sci. Paris 260(1965), 3831--3834.
\bibitem[51]{ch3:51} N. Th. Varopoulos, The functions that operate on $B_{0}(\Gamma)$ of a discrete group $\Gamma$, Bull. Soc. Math. France 93(1965), 301--321.
\bibitem[52]{ch3:52} N. Th. Varopoulos, Tensor algebras and harmonic analysis, Acta Math. 119(1967), 51--112.
\bibitem[53]{ch3:53} N. Wiener, Tauberian theorems, Ann. of Math. 33(1932), 1--100.
\bibitem[54]{ch3:54} N. Wiener and H. R. Pitt, On absolutely convergent Fourier-Stieltjes transforms, Duke Math. J. 4(1938), no. 2, 420--436.
\bibitem[55]{ch3:55} J. H. Williamson, A theorem on algebras of measures on topological groups, Proc. Edinburgh Math. Soc. 11(1958/1959), 195--206.
\bibitem[56]{ch3:56} K. Wozniakowski, A new proof of the restriction theorem for weak type (1,1) multipliers on $\mathbf{R}^n$, Illinois J. Math. 40(1996), 483.
\bibitem[57]{ch3:57} M. Zafran, On the spectra of multipliers, Pacific J. Math. 47(1973), 609--626.
\bibitem[58]{ch3:58} M. Zafran, The spectra of multiplier transformations on the $L_{p}$ spaces, Ann. of Math. 103(1976), no. 2, 355--374.
\bibitem[59]{ch3:59} M. Zafran, The functions operating on multiplier algebras, J. Functional Analysis 26(1977), no. 3, 289--314.
\bibitem[60]{ch3:60} M. Zafran, Exponential estimates in multiplier algebras, Conference on harmonic analysis in honor of Antoni Zygmund, Vol. I, 114--121, 1983.
\end{thebibliography}